\clearpage
\section{Appendix: Scriptural Date and Location}
% Generated. Do not edit.
% Deterministic projection of research/chronology.toml.
% schema: 3
% document: liturgy/roman-rite/postconciliar/roman-missal-third-edition-en-us-2011/propers/temporal/pc-s54-twenty-eighth-sunday-in-ordinary-time-year-a
% generated-by: tools/tpt proper-chronology annotations
\newcommand{\chronologyannotationclaim}[6]{#6}
\newcommand{\chronologyannotationcomparisonclaim}[8]{#8}
\newcommand{\chronologyannotationreach}[3]{}
\newcommand{\chronologyannotationgroup}[3]{#3}
\newcommand{\chronologyannotationcomparisongroup}[4]{#4}
\newcommand{\chronologyannotation}[1]{%
  \ifcsname triptychchronologyannotation@#1\endcsname
    \csname triptychchronologyannotation@#1\endcsname
  \else
    \PackageError{triptych}{No chronology annotation for #1}{}%
  \fi}
% appointment-dependency: guidance/liturgy/postconciliar-propers-registry.md sha256=ffe8a6d8b7d79c652bdc49b3d5f6d4dbb32e3506abc7ab0310e745c2b7306ad4
% appointment-dependency: src/gpt/liturgy/roman-rite/postconciliar/roman-missal-third-edition-en-us-2011/propers/registry/README.md sha256=65db141bfd1ec45b556069b256d53a2e65c945e1d7ec9df8ffb07dd76ee0a2bd
% appointment-dependency: src/gpt/liturgy/roman-rite/postconciliar/roman-missal-third-edition-en-us-2011/propers/registry/formula-dispositions.md sha256=c232dc6b4f3fa676c971ba081765471f366de1509dac0ddabca403f012f59879
% appointment-dependency: src/gpt/liturgy/roman-rite/postconciliar/roman-missal-third-edition-en-us-2011/propers/temporal/pc-s54-twenty-eighth-sunday-in-ordinary-time-year-a/instance/manifest.md sha256=684b6a565026a5227b4a1c62817e40a6b9fef727d07d56218ceee78d123422ef
% appointment-dependency: src/gpt/liturgy/roman-rite/postconciliar/roman-missal-third-edition-en-us-2011/propers/temporal/pc-s54-twenty-eighth-sunday-in-ordinary-time-year-a/propers/verified.md sha256=3fc1045aed7de765bb183ff9e3097f7d46f78bb2ba4164e3be84d30042b1208e
% appointment-dependency: src/gpt/liturgy/roman-rite/postconciliar/roman-missal-third-edition-en-us-2011/propers/temporal/pc-s54-twenty-eighth-sunday-in-ordinary-time-year-a/research/chronology-inputs.toml sha256=110e8e186301d2cfe4abece105f7028e7a6d1415445434971af4e7303df4332a
% appointment-dependency: src/gpt/liturgy/roman-rite/postconciliar/roman-missal-third-edition-en-us-2011/propers/temporal/shared/ordinary-time/weeks/28/propers/verified.md sha256=0affca6cbeaee7ee3111b2d8aaa5ad4696fafd00e82673fa32251987095ad657
% appointment-note: entrance: Psalm 129:3-4; identified-basis; vulgate numbering; canonical verse queries, not appointed wording.
% appointment-note: collect: Missal oration without an appointed biblical citation; thematic affinities do not create a Scripture appointment.
% appointment-note: first-reading: Isaiah 25:6-10a; appointed; vulgate numbering; whole-verse query envelope, not subverse chronology or appointed wording.
% appointment-note: psalm: Psalm 23:1-3a,3b-4,5,6; appointed; hebrew numbering; whole-verse query envelope, not subverse chronology or appointed wording.
% appointment-note: psalm: Psalm 23:6cd; response; hebrew numbering; whole-verse query envelope, not subverse chronology or appointed wording.
% appointment-note: second-reading: Philippians 4:12-14,19-20; appointed; vulgate numbering; canonical verse queries, not appointed wording.
% appointment-note: acclamation: Cf. Ephesians 1:17-18; adaptation; vulgate numbering; canonical verse queries, not appointed wording.
% appointment-note: gospel-long: Matthew 22:1-14; appointed; vulgate numbering; canonical verse queries, not appointed wording.
% appointment-note: gospel-short: Matthew 22:1-10; appointed; vulgate numbering; canonical verse queries, not appointed wording.
% appointment-note: offerings: Missal oration without an appointed biblical citation; thematic affinities do not create a Scripture appointment.
% appointment-note: communion-psalm: Cf. Psalm 33:11; adaptation; vulgate numbering; canonical verse queries, not appointed wording.
% appointment-note: communion-john: 1 John 3:2; identified-basis; vulgate numbering; canonical verse queries, not appointed wording.
% appointment-note: after-communion: Missal oration without an appointed biblical citation; thematic affinities do not create a Scripture appointment.
% comparison-dependency: src/gpt/liturgy/roman-rite/postconciliar/roman-missal-third-edition-en-us-2011/propers/temporal/pc-s54-twenty-eighth-sunday-in-ordinary-time-year-a/research/chronology-profile-comparisons.toml sha256=e77bdc39d80aba55b7dbae9c0d2426251d61ca8b0fb7209a94fc65293a23a82b
% entrance: Entrance antiphon
\expandafter\def\csname triptychchronologyannotation@entrance\endcsname{%
  \chronologyannotationgroup{historical-setting}{research-pending}{\textbf{Historical setting} -- Occasion date unresolved.}%
  \space \chronologyannotationgroup{composition}{preferred}{\textbf{Composition}: \chronologyannotationreach{Ps.129.3}{Ps}{true}\chronologyannotationreach{Ps.129.4}{Ps}{true}\chronologyannotationclaim{critical.psalms.latest-composition-boundary}{composition}{catholic-critical-v1}{preferred}{before the Maccabean period, around 165 B.C.; no individual psalm can be dated securely}{Before c. 165 B.C.}}%
}
% first-reading: First reading
\expandafter\def\csname triptychchronologyannotation@first-reading\endcsname{%
  \chronologyannotationgroup{traditional-attribution}{preferred}{\textbf{Traditional attribution}: \chronologyannotationreach{Is.25.10}{Is}{true}\chronologyannotationreach{Is.25.6}{Is}{true}\chronologyannotationreach{Is.25.7}{Is}{true}\chronologyannotationreach{Is.25.8}{Is}{true}\chronologyannotationreach{Is.25.9}{Is}{true}\chronologyannotationclaim{israel.prophets.isaias-traditional-era}{traditional-attribution}{catholic-traditional-v1}{preferred}{from the closing year of Ozias, King of Juda (740 B.C.); no prophecies appear to be later than 701 B.C.}{Isaias (ministry in Souvay's traditional account), B.C. 740--701}.}%
}
% psalm: Responsorial Psalm
\expandafter\def\csname triptychchronologyannotation@psalm\endcsname{%
  \chronologyannotationgroup{historical-setting}{research-pending}{\textbf{Historical setting} -- Occasion date unresolved.}%
  \space \chronologyannotationgroup{composition}{preferred}{\textbf{Composition}: \chronologyannotationreach{Ps.22.1}{Ps}{true}\chronologyannotationreach{Ps.22.2}{Ps}{true}\chronologyannotationreach{Ps.22.3}{Ps}{true}\chronologyannotationreach{Ps.22.4}{Ps}{true}\chronologyannotationreach{Ps.22.5}{Ps}{true}\chronologyannotationreach{Ps.22.6}{Ps}{true}\chronologyannotationclaim{critical.psalms.latest-composition-boundary}{composition}{catholic-critical-v1}{preferred}{before the Maccabean period, around 165 B.C.; no individual psalm can be dated securely}{Before c. 165 B.C.}}%
}
% second-reading: Second reading
\expandafter\def\csname triptychchronologyannotation@second-reading\endcsname{%
  \chronologyannotationgroup{composition}{disputed}{\textbf{Composition} -- disputed: \chronologyannotationreach{Phil.4.12}{Phil}{true}\chronologyannotationreach{Phil.4.13}{Phil}{true}\chronologyannotationreach{Phil.4.14}{Phil}{true}\chronologyannotationreach{Phil.4.19}{Phil}{true}\chronologyannotationreach{Phil.4.20}{Phil}{true}\chronologyannotationclaim{composition.epistle-to-the-philippians}{composition}{catholic-traditional-v1}{disputed}{(Philemon; Colossians; Ephesians; Philippians), 61}{A.D. 61}; \chronologyannotationreach{Phil.4.12}{Phil}{true}\chronologyannotationreach{Phil.4.13}{Phil}{true}\chronologyannotationreach{Phil.4.14}{Phil}{true}\chronologyannotationreach{Phil.4.19}{Phil}{true}\chronologyannotationreach{Phil.4.20}{Phil}{true}\chronologyannotationclaim{composition.epistle-to-the-philippians}{composition}{catholic-traditional-v1}{disputed}{at Rome (A.D. 62-64)}{A.D. 62--64}.}%
}
% acclamation: Gospel acclamation
\expandafter\def\csname triptychchronologyannotation@acclamation\endcsname{%
  \chronologyannotationgroup{composition}{disputed}{\textbf{Composition} -- disputed: \chronologyannotationreach{Eph.1.17}{Eph}{true}\chronologyannotationreach{Eph.1.18}{Eph}{true}\chronologyannotationclaim{composition.epistle-to-the-ephesians}{composition}{catholic-traditional-v1}{disputed}{a period between 58 and 63}{A.D. 58--63}; \chronologyannotationreach{Eph.1.17}{Eph}{true}\chronologyannotationreach{Eph.1.18}{Eph}{true}\chronologyannotationclaim{composition.epistle-to-the-ephesians}{composition}{catholic-traditional-v1}{disputed}{(Philemon; Colossians; Ephesians; Philippians), 61}{A.D. 61}.}%
}
% gospel-long: Gospel: long form
\expandafter\def\csname triptychchronologyannotation@gospel-long\endcsname{%
  \chronologyannotationgroup{narrated-event}{preferred}{\textbf{Event}: \chronologyannotationreach{Matt.22.1}{Matt.22.1-14}{false}\chronologyannotationreach{Matt.22.10}{Matt.22.1-14}{false}\chronologyannotationreach{Matt.22.11}{Matt.22.1-14}{false}\chronologyannotationreach{Matt.22.12}{Matt.22.1-14}{false}\chronologyannotationreach{Matt.22.13}{Matt.22.1-14}{false}\chronologyannotationreach{Matt.22.14}{Matt.22.1-14}{false}\chronologyannotationreach{Matt.22.2}{Matt.22.1-14}{false}\chronologyannotationreach{Matt.22.3}{Matt.22.1-14}{false}\chronologyannotationreach{Matt.22.4}{Matt.22.1-14}{false}\chronologyannotationreach{Matt.22.5}{Matt.22.1-14}{false}\chronologyannotationreach{Matt.22.6}{Matt.22.1-14}{false}\chronologyannotationreach{Matt.22.7}{Matt.22.1-14}{false}\chronologyannotationreach{Matt.22.8}{Matt.22.1-14}{false}\chronologyannotationreach{Matt.22.9}{Matt.22.1-14}{false}\chronologyannotationclaim{life-of-christ.parable-of-the-marriage-feast}{narrated-event}{catholic-traditional-v1}{preferred}{derived A.D. 29 from February, A.U.C. 782- Passover, 782}{The parable of the marriage feast, proposed in the Temple, A.D. 29 (derived)}.}%
  \space \chronologyannotationgroup{composition}{disputed}{\textbf{Composition} -- disputed: \chronologyannotationreach{Matt.22.1}{Matt}{true}\chronologyannotationreach{Matt.22.10}{Matt}{true}\chronologyannotationreach{Matt.22.11}{Matt}{true}\chronologyannotationreach{Matt.22.12}{Matt}{true}\chronologyannotationreach{Matt.22.13}{Matt}{true}\chronologyannotationreach{Matt.22.14}{Matt}{true}\chronologyannotationreach{Matt.22.2}{Matt}{true}\chronologyannotationreach{Matt.22.3}{Matt}{true}\chronologyannotationreach{Matt.22.4}{Matt}{true}\chronologyannotationreach{Matt.22.5}{Matt}{true}\chronologyannotationreach{Matt.22.6}{Matt}{true}\chronologyannotationreach{Matt.22.7}{Matt}{true}\chronologyannotationreach{Matt.22.8}{Matt}{true}\chronologyannotationreach{Matt.22.9}{Matt}{true}\chronologyannotationclaim{composition.gospel-of-matthew}{composition}{catholic-traditional-v1}{disputed}{about A.D. 38-45}{c. A.D. 38--45}; \chronologyannotationreach{Matt.22.1}{Matt}{true}\chronologyannotationreach{Matt.22.10}{Matt}{true}\chronologyannotationreach{Matt.22.11}{Matt}{true}\chronologyannotationreach{Matt.22.12}{Matt}{true}\chronologyannotationreach{Matt.22.13}{Matt}{true}\chronologyannotationreach{Matt.22.14}{Matt}{true}\chronologyannotationreach{Matt.22.2}{Matt}{true}\chronologyannotationreach{Matt.22.3}{Matt}{true}\chronologyannotationreach{Matt.22.4}{Matt}{true}\chronologyannotationreach{Matt.22.5}{Matt}{true}\chronologyannotationreach{Matt.22.6}{Matt}{true}\chronologyannotationreach{Matt.22.7}{Matt}{true}\chronologyannotationreach{Matt.22.8}{Matt}{true}\chronologyannotationreach{Matt.22.9}{Matt}{true}\chronologyannotationclaim{composition.gospel-of-matthew}{composition}{catholic-traditional-v1}{disputed}{about the year 40-42}{c. A.D. 40--42}; \chronologyannotationreach{Matt.22.1}{Matt}{true}\chronologyannotationreach{Matt.22.10}{Matt}{true}\chronologyannotationreach{Matt.22.11}{Matt}{true}\chronologyannotationreach{Matt.22.12}{Matt}{true}\chronologyannotationreach{Matt.22.13}{Matt}{true}\chronologyannotationreach{Matt.22.14}{Matt}{true}\chronologyannotationreach{Matt.22.2}{Matt}{true}\chronologyannotationreach{Matt.22.3}{Matt}{true}\chronologyannotationreach{Matt.22.4}{Matt}{true}\chronologyannotationreach{Matt.22.5}{Matt}{true}\chronologyannotationreach{Matt.22.6}{Matt}{true}\chronologyannotationreach{Matt.22.7}{Matt}{true}\chronologyannotationreach{Matt.22.8}{Matt}{true}\chronologyannotationreach{Matt.22.9}{Matt}{true}\chronologyannotationclaim{composition.gospel-of-matthew}{composition}{catholic-traditional-v1}{disputed}{the years 40-45}{A.D. 40--45}; \chronologyannotationreach{Matt.22.1}{Matt}{true}\chronologyannotationreach{Matt.22.10}{Matt}{true}\chronologyannotationreach{Matt.22.11}{Matt}{true}\chronologyannotationreach{Matt.22.12}{Matt}{true}\chronologyannotationreach{Matt.22.13}{Matt}{true}\chronologyannotationreach{Matt.22.14}{Matt}{true}\chronologyannotationreach{Matt.22.2}{Matt}{true}\chronologyannotationreach{Matt.22.3}{Matt}{true}\chronologyannotationreach{Matt.22.4}{Matt}{true}\chronologyannotationreach{Matt.22.5}{Matt}{true}\chronologyannotationreach{Matt.22.6}{Matt}{true}\chronologyannotationreach{Matt.22.7}{Matt}{true}\chronologyannotationreach{Matt.22.8}{Matt}{true}\chronologyannotationreach{Matt.22.9}{Matt}{true}\chronologyannotationclaim{composition.gospel-of-matthew}{composition}{catholic-traditional-v1}{disputed}{about the year 60-68}{c. A.D. 60--68}; \chronologyannotationreach{Matt.22.1}{Matt}{true}\chronologyannotationreach{Matt.22.10}{Matt}{true}\chronologyannotationreach{Matt.22.11}{Matt}{true}\chronologyannotationreach{Matt.22.12}{Matt}{true}\chronologyannotationreach{Matt.22.13}{Matt}{true}\chronologyannotationreach{Matt.22.14}{Matt}{true}\chronologyannotationreach{Matt.22.2}{Matt}{true}\chronologyannotationreach{Matt.22.3}{Matt}{true}\chronologyannotationreach{Matt.22.4}{Matt}{true}\chronologyannotationreach{Matt.22.5}{Matt}{true}\chronologyannotationreach{Matt.22.6}{Matt}{true}\chronologyannotationreach{Matt.22.7}{Matt}{true}\chronologyannotationreach{Matt.22.8}{Matt}{true}\chronologyannotationreach{Matt.22.9}{Matt}{true}\chronologyannotationclaim{composition.gospel-of-matthew}{composition}{catholic-traditional-v1}{disputed}{about the years 64-67}{c. A.D. 64--67}; \chronologyannotationreach{Matt.22.1}{Matt}{true}\chronologyannotationreach{Matt.22.10}{Matt}{true}\chronologyannotationreach{Matt.22.11}{Matt}{true}\chronologyannotationreach{Matt.22.12}{Matt}{true}\chronologyannotationreach{Matt.22.13}{Matt}{true}\chronologyannotationreach{Matt.22.14}{Matt}{true}\chronologyannotationreach{Matt.22.2}{Matt}{true}\chronologyannotationreach{Matt.22.3}{Matt}{true}\chronologyannotationreach{Matt.22.4}{Matt}{true}\chronologyannotationreach{Matt.22.5}{Matt}{true}\chronologyannotationreach{Matt.22.6}{Matt}{true}\chronologyannotationreach{Matt.22.7}{Matt}{true}\chronologyannotationreach{Matt.22.8}{Matt}{true}\chronologyannotationreach{Matt.22.9}{Matt}{true}\chronologyannotationclaim{composition.gospel-of-matthew}{composition}{catholic-traditional-v1}{disputed}{about the year 50}{c. A.D. 50}.}%
  \space \chronologyannotationcomparisongroup{catholic-critical-v1}{composition}{composition-only}{\textbf{Composition (Catholic critical chronology, for comparison)}: \chronologyannotationreach{Matt.22.1}{Matt}{true}\chronologyannotationreach{Matt.22.10}{Matt}{true}\chronologyannotationreach{Matt.22.11}{Matt}{true}\chronologyannotationreach{Matt.22.12}{Matt}{true}\chronologyannotationreach{Matt.22.13}{Matt}{true}\chronologyannotationreach{Matt.22.14}{Matt}{true}\chronologyannotationreach{Matt.22.2}{Matt}{true}\chronologyannotationreach{Matt.22.3}{Matt}{true}\chronologyannotationreach{Matt.22.4}{Matt}{true}\chronologyannotationreach{Matt.22.5}{Matt}{true}\chronologyannotationreach{Matt.22.6}{Matt}{true}\chronologyannotationreach{Matt.22.7}{Matt}{true}\chronologyannotationreach{Matt.22.8}{Matt}{true}\chronologyannotationreach{Matt.22.9}{Matt}{true}\chronologyannotationcomparisonclaim{critical.gospel-of-matthew}{composition}{catholic-critical-v1}{catholic-critical-v1}{preferred}{passage.united-states-conference-of-catholic-bishops.new-american-bible-revised-edition.english-usccb-web-2026-09-21.matthew-introduction}{post-A.D. 70 date}{The Greek Gospel of Matthew in the NABRE introduction, post-A.D. 70 date}.}%
}
% gospel-short: Gospel: short form
\expandafter\def\csname triptychchronologyannotation@gospel-short\endcsname{%
  \chronologyannotationgroup{narrated-event}{preferred}{\textbf{Event}: \chronologyannotationreach{Matt.22.1}{Matt.22.1-14}{false}\chronologyannotationreach{Matt.22.10}{Matt.22.1-14}{false}\chronologyannotationreach{Matt.22.2}{Matt.22.1-14}{false}\chronologyannotationreach{Matt.22.3}{Matt.22.1-14}{false}\chronologyannotationreach{Matt.22.4}{Matt.22.1-14}{false}\chronologyannotationreach{Matt.22.5}{Matt.22.1-14}{false}\chronologyannotationreach{Matt.22.6}{Matt.22.1-14}{false}\chronologyannotationreach{Matt.22.7}{Matt.22.1-14}{false}\chronologyannotationreach{Matt.22.8}{Matt.22.1-14}{false}\chronologyannotationreach{Matt.22.9}{Matt.22.1-14}{false}\chronologyannotationclaim{life-of-christ.parable-of-the-marriage-feast}{narrated-event}{catholic-traditional-v1}{preferred}{derived A.D. 29 from February, A.U.C. 782- Passover, 782}{The parable of the marriage feast, proposed in the Temple, A.D. 29 (derived)}.}%
  \space \chronologyannotationgroup{composition}{disputed}{\textbf{Composition} -- disputed: \chronologyannotationreach{Matt.22.1}{Matt}{true}\chronologyannotationreach{Matt.22.10}{Matt}{true}\chronologyannotationreach{Matt.22.2}{Matt}{true}\chronologyannotationreach{Matt.22.3}{Matt}{true}\chronologyannotationreach{Matt.22.4}{Matt}{true}\chronologyannotationreach{Matt.22.5}{Matt}{true}\chronologyannotationreach{Matt.22.6}{Matt}{true}\chronologyannotationreach{Matt.22.7}{Matt}{true}\chronologyannotationreach{Matt.22.8}{Matt}{true}\chronologyannotationreach{Matt.22.9}{Matt}{true}\chronologyannotationclaim{composition.gospel-of-matthew}{composition}{catholic-traditional-v1}{disputed}{about A.D. 38-45}{c. A.D. 38--45}; \chronologyannotationreach{Matt.22.1}{Matt}{true}\chronologyannotationreach{Matt.22.10}{Matt}{true}\chronologyannotationreach{Matt.22.2}{Matt}{true}\chronologyannotationreach{Matt.22.3}{Matt}{true}\chronologyannotationreach{Matt.22.4}{Matt}{true}\chronologyannotationreach{Matt.22.5}{Matt}{true}\chronologyannotationreach{Matt.22.6}{Matt}{true}\chronologyannotationreach{Matt.22.7}{Matt}{true}\chronologyannotationreach{Matt.22.8}{Matt}{true}\chronologyannotationreach{Matt.22.9}{Matt}{true}\chronologyannotationclaim{composition.gospel-of-matthew}{composition}{catholic-traditional-v1}{disputed}{about the year 40-42}{c. A.D. 40--42}; \chronologyannotationreach{Matt.22.1}{Matt}{true}\chronologyannotationreach{Matt.22.10}{Matt}{true}\chronologyannotationreach{Matt.22.2}{Matt}{true}\chronologyannotationreach{Matt.22.3}{Matt}{true}\chronologyannotationreach{Matt.22.4}{Matt}{true}\chronologyannotationreach{Matt.22.5}{Matt}{true}\chronologyannotationreach{Matt.22.6}{Matt}{true}\chronologyannotationreach{Matt.22.7}{Matt}{true}\chronologyannotationreach{Matt.22.8}{Matt}{true}\chronologyannotationreach{Matt.22.9}{Matt}{true}\chronologyannotationclaim{composition.gospel-of-matthew}{composition}{catholic-traditional-v1}{disputed}{the years 40-45}{A.D. 40--45}; \chronologyannotationreach{Matt.22.1}{Matt}{true}\chronologyannotationreach{Matt.22.10}{Matt}{true}\chronologyannotationreach{Matt.22.2}{Matt}{true}\chronologyannotationreach{Matt.22.3}{Matt}{true}\chronologyannotationreach{Matt.22.4}{Matt}{true}\chronologyannotationreach{Matt.22.5}{Matt}{true}\chronologyannotationreach{Matt.22.6}{Matt}{true}\chronologyannotationreach{Matt.22.7}{Matt}{true}\chronologyannotationreach{Matt.22.8}{Matt}{true}\chronologyannotationreach{Matt.22.9}{Matt}{true}\chronologyannotationclaim{composition.gospel-of-matthew}{composition}{catholic-traditional-v1}{disputed}{about the year 60-68}{c. A.D. 60--68}; \chronologyannotationreach{Matt.22.1}{Matt}{true}\chronologyannotationreach{Matt.22.10}{Matt}{true}\chronologyannotationreach{Matt.22.2}{Matt}{true}\chronologyannotationreach{Matt.22.3}{Matt}{true}\chronologyannotationreach{Matt.22.4}{Matt}{true}\chronologyannotationreach{Matt.22.5}{Matt}{true}\chronologyannotationreach{Matt.22.6}{Matt}{true}\chronologyannotationreach{Matt.22.7}{Matt}{true}\chronologyannotationreach{Matt.22.8}{Matt}{true}\chronologyannotationreach{Matt.22.9}{Matt}{true}\chronologyannotationclaim{composition.gospel-of-matthew}{composition}{catholic-traditional-v1}{disputed}{about the years 64-67}{c. A.D. 64--67}; \chronologyannotationreach{Matt.22.1}{Matt}{true}\chronologyannotationreach{Matt.22.10}{Matt}{true}\chronologyannotationreach{Matt.22.2}{Matt}{true}\chronologyannotationreach{Matt.22.3}{Matt}{true}\chronologyannotationreach{Matt.22.4}{Matt}{true}\chronologyannotationreach{Matt.22.5}{Matt}{true}\chronologyannotationreach{Matt.22.6}{Matt}{true}\chronologyannotationreach{Matt.22.7}{Matt}{true}\chronologyannotationreach{Matt.22.8}{Matt}{true}\chronologyannotationreach{Matt.22.9}{Matt}{true}\chronologyannotationclaim{composition.gospel-of-matthew}{composition}{catholic-traditional-v1}{disputed}{about the year 50}{c. A.D. 50}.}%
  \space \chronologyannotationcomparisongroup{catholic-critical-v1}{composition}{composition-only}{\textbf{Composition (Catholic critical chronology, for comparison)}: \chronologyannotationreach{Matt.22.1}{Matt}{true}\chronologyannotationreach{Matt.22.10}{Matt}{true}\chronologyannotationreach{Matt.22.2}{Matt}{true}\chronologyannotationreach{Matt.22.3}{Matt}{true}\chronologyannotationreach{Matt.22.4}{Matt}{true}\chronologyannotationreach{Matt.22.5}{Matt}{true}\chronologyannotationreach{Matt.22.6}{Matt}{true}\chronologyannotationreach{Matt.22.7}{Matt}{true}\chronologyannotationreach{Matt.22.8}{Matt}{true}\chronologyannotationreach{Matt.22.9}{Matt}{true}\chronologyannotationcomparisonclaim{critical.gospel-of-matthew}{composition}{catholic-critical-v1}{catholic-critical-v1}{preferred}{passage.united-states-conference-of-catholic-bishops.new-american-bible-revised-edition.english-usccb-web-2026-09-21.matthew-introduction}{post-A.D. 70 date}{The Greek Gospel of Matthew in the NABRE introduction, post-A.D. 70 date}.}%
}
% communion-psalm: Communion: psalm option
\expandafter\def\csname triptychchronologyannotation@communion-psalm\endcsname{%
  \chronologyannotationgroup{traditional-attribution}{disputed}{\textbf{Traditional attribution} -- disputed: \chronologyannotationreach{Ps.33.11}{Ps.33 Ps.39 Ps.70 Ps.85 Ps.94 Ps.97}{true}\chronologyannotationclaim{israel.monarchy.david-traditional-era}{traditional-attribution}{catholic-traditional-v1}{disputed}{reigned from 1055 to 1015 B.C.}{David (reign in the usual chronology), B.C. 1055--1015}.}%
  \space \chronologyannotationgroup{superscription-setting}{preferred}{\textbf{Superscription setting}: \chronologyannotationreach{Ps.33.11}{Ps.33 Ps.55}{true}\chronologyannotationclaim{israel.monarchy.david-at-geth}{superscription-setting}{catholic-traditional-v1}{preferred}{A. M. 2944, A. C. 1060}{David's flight to Achis, king of Geth, A.M. 2944}.}%
  \space \chronologyannotationgroup{composition}{preferred}{\textbf{Composition}: \chronologyannotationreach{Ps.33.11}{Ps}{true}\chronologyannotationclaim{critical.psalms.latest-composition-boundary}{composition}{catholic-critical-v1}{preferred}{before the Maccabean period, around 165 B.C.; no individual psalm can be dated securely}{Before c. 165 B.C.}}%
}
% communion-john: Communion: Johannine option
\expandafter\def\csname triptychchronologyannotation@communion-john\endcsname{%
  \chronologyannotationgroup{composition}{preferred}{\textbf{Composition}: \chronologyannotationreach{1John.3.2}{1John}{true}\chronologyannotationclaim{composition.first-epistle-of-st-john}{composition}{catholic-traditional-v1}{preferred}{from the year 90 to 100 (approximately)}{c. A.D. 90--100}.}%
}

\begin{dossiertable}
Psalm & Ps 22:1--6 & Pasture, valley and house are poetic scenes; no particular valley is established. & \chronodate{psalm}{\chronologyannotation{psalm}}\\
\dossierprose{The Douay title attributes the psalm to David. That attribution supplies no dated occasion. The broad composition limit belongs to the Psalter's historical horizon, not an exact date for this poem (NABRE, Psalms introduction).}
\midrule
Communion option & Cf. Ps 33:11 & The superscription's narrative setting is Geth (Gath); not an identified composition site. & \chronodate{communion-psalm}{\chronologyannotation{communion-psalm}}\\
\cmidrule(lr){2-4}
\dossierevent{Narrated setting: David's flight and assumed madness; the superscription names Achimelech, while the linked 1 Kings 21 narrative names Achis. The displayed A.M. frame is Haydock's narrative chronology, not a newly calculated date for the psalm.}
\dossierprose{The reign range identifies the traditionally attributed person; the flight precedes his reign. Authorship, superscription setting and composition remain different claims (Douay Ps 33 title; Corbett, ``David, King''; Haydock on 1 Kings chapters 20 through 22; NABRE, Psalms introduction).}
\midrule
Entrance & Ps 129:3--4 & The depths express distress and guilt, not an identified site. & \chronodate{entrance}{\chronologyannotation{entrance}}\\
\dossierprose{No author is named in the inspected Douay text. No particular historical occasion is established. The displayed composition limit is broad; it does not date an individual act of penitence (NABRE, Psalms introduction).}
\midrule
First reading & Is 25:6--10a & Zion/Jerusalem supplies ``this mountain'' after Is 24:23; the feast embraces all peoples. & \chronodate{first-reading}{\chronologyannotation{first-reading}}\\
\dossierprose{The displayed era concerns Isaiah's ministry in Souvay's traditional account. It is neither a composition date for chapter 25 nor a date for the promised feast. Souvay describes the prophetic referent as an undetermined future (``Isaias,'' Life and First Isaias).}
\midrule
Gospel, both forms & Mt 22:1--14; short 1--10 & Jerusalem Temple discourse; a proposed Antioch composition setting belongs to another question. & \chronodate{gospel-long}{\chronologyannotation{gospel-long}}\\
\cmidrule(lr){2-4}
\dossierevent{Narrated event: Jesus' teaching in the Temple after Mt 21:23. The displayed event date is derived within Maas's traditional chronology; it does not date a historical royal wedding. Both Gospel forms share this event and composition dossier.}
\dossierprose{The traditional composition proposals are disputed historical hypotheses, not six equally certain findings. Jacquier calls the early dispersal tradition unreliable; his later reckoning is conditional and the Irenaeus inference inconclusive. Durand's earlier dating concerns an Aramaic predecessor, with its Greek rendering undated. The separately displayed Catholic-critical comparison concerns the Greek Gospel: the NABRE introduction favors a date after Jerusalem's destruction, probably in the following decade, and treats Antioch as a plausible place. These are distinct accounts, not a merged range.}
\end{dossiertable}
\clearpage
\begin{dossiertable}
Acclamation & Cf. Eph 1:17--18 & Ephesus / wider recipients; Caesarea or Rome proposed for writing. & \chronodate{acclamation}{\chronologyannotation{acclamation}}\\
\dossierprose{The displayed composition alternatives come from the historical discussions of Ladeuze and Prat. Addressees and writing place differ; the adapted liturgical verse receives no separate biblical event date.}
\midrule
Second reading & Phil 4:12--14, 19--20 & Philippi is the addressee's city; the cited traditional account places writing in Rome. & \chronodate{second-reading}{\chronologyannotation{second-reading}}\\
\dossierprose{Prat and Vander Heeren preserve alternative Roman-captivity datings. Their answers are historical judgments, not a single composite range or an assertion of modern consensus. The particular hunger and gifts are not independently dated incidents.}
\midrule
Communion option & 1 Jn 3:2 & No exact writing site established here; future sight has no earthly location. & \chronodate{communion-john}{\chronologyannotation{communion-john}}\\
\dossierprose{The displayed range concerns composition, not the future appearing or vision of God (Durand, ``The New Testament''; Drum, ``Epistles of Saint John''). Present adoption and final manifestation are theological relations within the passage, not two dated earthly events.}
\end{dossiertable}
The three orations are composed liturgical prayers and have no directly appointed biblical locus to date. Psalm numbering above is Vulgate; Hebrew equivalents appear in the opening inventory. The chronological proposals retain their particular relation to authorship, setting, composition or narrated event. They do not fix the date of a promise's fulfillment.
